% Methodology: problem definition and the evaluation protocol (issue #7).
% Features are in feature_engineering.tex and models in models.tex.

\section{Methodology}
\label{sec:methodology}

This section defines the prediction task and the evaluation protocol shared by
every Phase~1 model; Phase~2 keeps it unchanged. Fig.~\ref{fig:protocol}
summarizes the protocol. The features and models it evaluates are described in
Sections~\ref{sec:feature_engineering} and~\ref{sec:models}.

\begin{figure}[!t]
\centering
\begin{tikzpicture}[
    node distance=3.2mm,
    box/.style={draw, rounded corners=2pt, align=center, font=\footnotesize,
                 text width=0.74\columnwidth, inner sep=3pt},
    arrow/.style={-{Stealth[length=2mm]}, thick},
    note/.style={font=\scriptsize\itshape, align=left}]
  \node[box] (train) {100 FD001 training engines, run to failure;\\
    target $y_{u,t}=T_u-t$ (uncapped)};
  \node[box, below=of train] (folds) {3 engine-grouped folds
    (\texttt{GroupKFold} on \texttt{unit}, seed 42)};
  \node[box, below=of folds, fill=black!6] (fit) {Fit on the training fold only:
    select sensors $\rightarrow$ regime z-score $\rightarrow$ rolling features
    $\rightarrow$ model};
  \node[box, below=of fit] (cv) {Predict every cycle of the held-out engines;
    choose the grid candidate with the lowest mean CV RMSE};
  \node[box, below=of cv] (refit) {Refit the chosen pipeline on all 100 training
    engines};
  \node[box, below=of refit] (test) {Predict every recorded cycle of the 100 test
    engines; score each once at its last cycle against \texttt{RUL\_FD001.txt}};
  \draw[arrow] (train) -- (folds);
  \draw[arrow] (folds) -- (fit);
  \draw[arrow] (fit) -- (cv);
  \draw[arrow] (cv) -- (refit);
  \draw[arrow] (refit) -- (test);
  \draw[arrow, dashed] (cv.east) -- ++(2.5mm,0) |- (fit.east)
    node[note, pos=0.25, right] {each fold,\\each candidate};
\end{tikzpicture}
\caption{Evaluation protocol. Every data-dependent preprocessing step is part of
the pipeline (shaded), so it is refitted with the model inside each training
fold.}
\label{fig:protocol}
\end{figure}

\subsection{Problem Formulation}
Engine $u$ is observed at cycles $t=1,\dots,T_u$ with predictor vector
$\mathbf{x}_{u,t}$. A training engine runs until it fails after cycle $T_u$, so
its RUL at cycle $t$ is $y_{u,t}=T_u-t$. A test engine is observed only up to a
cycle $\tau_u$ before failure, and its true RUL is given for that cycle alone.
A predictor may use only the engine's own history,
$\hat{y}_{u,t}=f(\mathbf{x}_{u,1},\dots,\mathbf{x}_{u,t})$, which matches how RUL
would be estimated during operation. The project brief asks for the target to be
capped only when stated, so the uncapped $y_{u,t}$ is the primary target. A
target capped at 125 cycles, $\min(y_{u,t},125)$, which much of the literature
uses~\cite{Heimes2008recurrent,Zheng2017long}, is treated only as an ablation.

\subsection{Engine-Grouped Validation}
Consecutive cycles of one engine are nearly identical, so splitting rows at
random would put cycles of the same engine on both sides of a split and
overstate accuracy. All validation therefore splits by engine: the 100 training
engines are divided into three folds with \texttt{GroupKFold} on the engine
identifier and a fixed seed of 42, and every engine is held out exactly once.
Cross-validation scores every cycle of the held-out engines, whose RUL ranges up
to 361 cycles, whereas the test set is scored at last observed cycles whose RUL
is between 7 and 145. Cross-validation error is therefore used only to choose
hyperparameters, not as an estimate of test error.

\subsection{Leakage Control}
Every step that learns from data is fitted on training rows only. Sensor
selection, the operating-regime clusters and their per-regime statistics are
part of a scikit-learn pipeline whose final step is the model, so a grid search
refits them inside each training fold, and the final model refits them on all
training engines before transforming the test engines. Rolling features use
only the current and earlier cycles of the same engine, so they are computed
the same way for a complete training trajectory and for a truncated test
engine. Unit tests check these properties: folds never share an engine, rolling
features ignore later cycles, test rows are scaled with training statistics,
and the model never receives the true RUL.

\subsection{Feature Sets and Tuning}
Each model is trained on two inputs. The raw set contains the cycle, the three
operating settings and all 21 sensors (25 predictors). The engineered set is
the 155 predictors of Section~\ref{sec:feature_engineering}. The engineered
set is the Phase~1 method, fixed before the test set was scored, and the raw set
is the comparison that measures what feature engineering adds. Each model is
tuned by grid search over the candidates in Section~\ref{sec:models}, selected
by mean CV RMSE, and the best candidate is refitted on all 100 training engines.

\subsection{Official Test Evaluation}
The fitted pipeline predicts every recorded cycle of each test engine from that
engine's complete history, so that the rolling features at the last cycle have
their past. Each engine is then scored once, at its last cycle, against the
true RUL in \texttt{RUL\_FD001.txt}. All 100 test engines are scored for every
model, and the per-engine predictions are saved with the metrics so that later
phases can compare models engine by engine.

\subsection{Metrics}
With $d_u=\hat{y}_u-y_u$ the error at the last cycle of test engine $u$, we
report the RMSE, the MAE and the asymmetric score of the PHM08
challenge~\cite{Saxena2008damage},
\begin{equation}
    S=\sum_{u}
    \begin{cases}
        e^{-d_u/13}-1, & d_u<0,\\
        e^{d_u/10}-1, & d_u\ge 0,
    \end{cases}
\end{equation}
which penalizes late predictions ($d_u>0$, more life predicted than remains)
more than early ones, because an unexpected failure costs more than early
maintenance. $S$ is a sum, not a mean, and grows exponentially, so one large
overestimate can dominate it; we always read it together with RMSE and MAE.
Following the health-stage view of prognostics~\cite{Lei2018machinery}, errors
are also reported by degradation stage of the true RUL: late
($y\le 50$), mid ($50<y\le 100$) and early ($y>100$). The FD001 test set has 33,
34 and 33 engines in these stages.

\subsection{Implementation and Reproducibility}
All shared code is in a Python package (\texttt{rul}) with 99 unit tests.
Notebooks 01 to 04 run from the raw NASA files to the saved metrics, figures and
per-engine predictions. Every stochastic step uses seed 42, and repeated runs
reproduce the saved results exactly.
